\documentclass[dvipdfmx,11pt,a4paper]{article}
\usepackage{amsmath,amssymb,amsthm,amsfonts}
\usepackage{geometry}
\geometry{margin=1.1in}
\usepackage{enumitem}
\usepackage{booktabs}
\usepackage{mathtools}
\usepackage[dvipdfmx]{hyperref}
\usepackage[nameinlink,noabbrev]{cleveref}

\newtheorem{theorem}{Theorem}[section]
\newtheorem{proposition}[theorem]{Proposition}
\newtheorem{lemma}[theorem]{Lemma}
\newtheorem{corollary}[theorem]{Corollary}
\theoremstyle{definition}
\newtheorem{definition}[theorem]{Definition}
\newtheorem{example}[theorem]{Example}
\theoremstyle{remark}
\newtheorem{remark}[theorem]{Remark}

\newcommand{\Jnorm}{J_{\rm norm}}
\newcommand{\Jbound}{\frac{3\pi^2}{8}}

\title{Pseudo-Quantum Dual-Spacetime 4-Valued Logic (PQ4L):\\
Amplitudes and Geometric Interference}

\author{https://github.com/hypernumbernet}
\date{\today}

\begin{document}
\maketitle

\begin{abstract}
We fix a named formal system, PQ4L, on top of the already-checked Dual
Spacetime 4-valued logic (D4L). A proposition's primary carrier is an
admissible torsion configuration, equivalently its PGA bivector
\(\Omega\) and rotor \(R=\exp\Omega\). The four labels
\(\lvert T\rangle,\lvert U\rangle,\lvert F\rangle,\lvert B\rangle\)
appear only after measuring the normalised height \(\Jnorm\).
Commutative overlap is the parameter Killing form; non-commutative
interference is the geometric commutator of torsion bivectors. Two
preorders on \([-1,1]\) separate height from saturation. None of this
is Birkhoff--von Neumann quantum logic. Hilbert space, the Born rule,
and an orthomodular lattice are not claimed. The statements listed as
theorems below are machine-checked in Lean~4.
\end{abstract}

\section{Position}

D4L \cite{D4L2026} classifies admissible configurations by the sign and
saturation of \(\Jnorm\in[-1,1]\) and equips the interval with
\(\min/\max/-\). Those connectives are distributive. They cannot see
the geometric product already present in the PGA core \cite{Diophantine2026}.

PQ4L does not replace that scalar layer. It raises the carrier one
step, to the configuration itself, and reads DST operations as logical
ones. The analogy is the same as pseudo-quantum methods in annealing:
there is an amplitude, a pairing, an interference term, and a
collapse, but there is no \(\hbar\) and no claim to orthomodularity.

The Phase~1 definitions \(\texttt{classifyOfMem}\), \(\min/\max\), and
the written potential \(U\) are left untouched. The module is a
parallel track: it is not re-exported from \(\texttt{DstDiophantine.Basic}\).

\section{Amplitudes}

\begin{definition}[Amplitude]
\label{def:amp}
An \emph{amplitude} is an admissible continuous torsion configuration
\(p=(\alpha_a,\beta_a)\). Write
\[
\Omega(p)=\texttt{omegaTorsion}\,p,\qquad
R(p)=\texttt{rotorTorsion}\,p=\exp\Omega(p).
\]
Measurement and collapse are the existing maps
\[
\langle p\rangle:=\Jnorm(p),\qquad
\mathrm{collapse}(p):=\texttt{ofParams}\,p\in\{\lvert T\rangle,\lvert U\rangle,\lvert F\rangle,\lvert B\rangle\}.
\]
The adjoint is the usual--dual swap \(p^\dagger=\texttt{daggerParams}\,p\).
\end{definition}

\begin{theorem}[Adjoint]
\label{thm:adjoint}
The swap is an involution and flips the observable:
\( (p^\dagger)^\dagger=p \) and \(\langle p^\dagger\rangle=-\langle p\rangle\).
Every D4L label is realised by some amplitude.
\end{theorem}

The four names are still not a function of each other under adjoint:
deepest \(\lvert B\rangle\) maps to \(\lvert F\rangle\), interior
\(\lvert B\rangle\) maps to \(\lvert U\rangle\) \cite{D4L2026}. Scalar
conjunction remains non-explosive: \(j\wedge\neg j=-|j|\) never
saturates \(\lvert F\rangle\).

\section{Two orders, not Belnap FOUR}

\begin{definition}
On \(\mathbb{R}\), and then by restriction on \([-1,1]\),
\[
j\preceq_\tau k \iff j\le k,
\qquad
j\preceq_\iota k \iff \lvert j\rvert\le\lvert k\rvert.
\]
The first is the height (boost-versus-rotation) order. The second is
the information (saturation) order.
\end{definition}

\begin{theorem}[Information bottom and tops]
\label{thm:info}
On \([-1,1]\), \(j=0\) is the unique information bottom, and the
information tops are exactly the walls \(j=\pm 1\). Equivalently, the
information bottom is the label \(\lvert T\rangle\), and the two tops
are \(\lvert F\rangle\) and deepest \(\lvert B\rangle\).
\end{theorem}

Interior \(\lvert B\rangle\) is strictly below deepest \(\lvert B\rangle\)
in the information order. The labels do not match Belnap's
\(\{\bot,t,f,\top\}\); no bilattice isomorphism is claimed.

\section{Geometric operations}

\begin{definition}[Overlap and interference]
\[
\langle p\mid q\rangle := \texttt{killingForm}\,p\,q
= 8\sum_{a=1}^3\bigl(\alpha_a\alpha'_a-\beta_a\beta'_a\bigr),
\qquad
[\Omega(p),\Omega(q)] := \Omega(p)\Omega(q)-\Omega(q)\Omega(p).
\]
Rotor composition \(R(p)R(q)\) is a product in the PGA. It is not
pulled back through the Baker--Campbell--Hausdorff formula to a third
torsion configuration.
\end{definition}

\begin{theorem}[Pairing and non-commutativity]
\label{thm:geo}
The overlap is symmetric, and \(\langle p\mid p\rangle=16\,J(p)\),
equivalently a constant multiple of \(\Jnorm(p)\). Distinct-axis
torsion bivectors need not commute: if \(p\) is a pure boost on axis
\(0\) and \(q\) a pure rotation on axis \(1\), then
\([\Omega(p),\Omega(q)]\ne 0\). Same-axis hyperbolic and cyclic
generators do commute, so the corresponding one-axis interference
vanishes.
\end{theorem}

The overlap is not a Born probability. Sandwich evaluation
\(R(p)\,v\,R(p)^\sim\) is restated independently of the gravity chart
layer; preservation of the degenerate quadratic form is not claimed.

The scalar connectives remain \(\neg j=-j\), \(j\wedge k=\min(j,k)\),
\(j\vee k=\max(j,k)\). They are distributive and non-explosive
\cite{D4L2026}. Geometric composition is associative as a PGA product
and is not claimed to be a lattice operation.

\section{What PQ4L does and does not show}

\begin{itemize}[leftmargin=1.4em]
  \item It \emph{does} make the PGA torsion bivector and rotor the
        carriers of a proposition, with \(\Jnorm\) as the observable.
  \item It \emph{does} exhibit a commutative pairing and a
        non-commutative interference term already proved in the
        algebra.
  \item It \emph{does} separate height from saturation by two
        preorders on \([-1,1]\).
  \item It does \emph{not} construct a Hilbert space, a Born rule, or
        an orthomodular lattice.
  \item It does \emph{not} identify the idea-note scale \(\lambda\)
        with CGA dilation or with \(\texttt{scaleTorsion}\).
  \item It does \emph{not} obtain a four-to-two collapse from the
        written potential \(U\), and it does \emph{not} refute
        G\"odel's theorems in classical arithmetic.
  \item It does \emph{not} pull \(R(p)R(q)\) back to a torsion
        configuration when \([\Omega(p),\Omega(q)]\ne 0\).
\end{itemize}

The Lean entry is \(\texttt{DstDiophantine.Logic}\). The PQ4L modules
are \(\texttt{Amplitude}\), \(\texttt{Order}\), and \(\texttt{Geometric}\).

\begin{thebibliography}{9}
\raggedright

\bibitem{D4L2026}
Dual Spacetime 4-Valued Paraconsistent Logic (D4L),
\texttt{dst-4-valued-logic.tex} (2026).

\bibitem{Diophantine2026}
Discrete Biquaternionic Double Spacetime,
\texttt{dst-diophantine.tex} (2026). Lean companion:
\url{https://github.com/hypernumbernet/DstDiophantine}.

\bibitem{DST2025}
Dual Spacetime Theory: Gravity as Torsion between Dual Spacetimes (2025).

\end{thebibliography}

\end{document}
